\documentclass[10pt,a4paper]{article}
\usepackage[margin=0.62in]{geometry}
\usepackage[T1]{fontenc}
\usepackage[default]{sourcesanspro}
\usepackage{microtype}
\usepackage{enumitem}
\usepackage{titlesec}
\usepackage{xcolor}
\usepackage[hidelinks]{hyperref}
\usepackage{tabularx}
\pagestyle{empty}

\input{glyphtounicode}
\pdfgentounicode=1

\definecolor{accent}{HTML}{176E87}
\definecolor{secondary}{HTML}{5F686F}
\definecolor{rulegray}{HTML}{CDD4D8}

\titleformat{\section}
  {\large\bfseries\color{accent}}
  {}{0pt}{}
  [\vspace{-0.42em}\color{rulegray}\rule{\linewidth}{0.4pt}]
\titlespacing*{\section}{0pt}{0.62em}{0.28em}
\setlist[itemize]{leftmargin=1.15em,itemsep=0.12em,topsep=0.12em,parsep=0pt,partopsep=0pt}
\setlength{\parindent}{0pt}
\setlength{\parskip}{0pt}
\setlength{\emergencystretch}{1em}
\hypersetup{
  pdftitle={Qi Ding — Academic Curriculum Vitae},
  pdfauthor={Qi Ding},
  pdfsubject={End-use transitions, spatial infrastructure, and power-system integration},
  pdfkeywords={energy systems, transport decarbonization, hydrogen, power systems, optimization}
}

\newcommand{\entry}[3]{%
  \begin{tabularx}{\linewidth}{@{}Xr@{}}
    \textbf{#1} & \textcolor{secondary}{#3} \\
    #2 &
  \end{tabularx}\vspace{0.04em}}

\newcommand{\project}[4]{%
  \item \begin{tabularx}{\linewidth}{@{}Xr@{}}
    \textbf{#1 --- #2} & \textcolor{secondary}{#4}
  \end{tabularx}\vspace{-0.32em}\\
  \textit{#3}}

\newcommand{\currentitem}[3]{%
  \item \textbf{#1} \hfill \textcolor{secondary}{#2}\\[-0.2em]
  #3}

\newcommand{\experience}[4]{%
  \begin{tabularx}{\linewidth}{@{}Xr@{}}
    \textbf{#1} --- #2 & \textcolor{secondary}{#4} \\
    \multicolumn{2}{@{}p{\dimexpr\linewidth\relax}@{}}{#3}
  \end{tabularx}\vspace{0.12em}}

\begin{document}

\begin{center}
  {\LARGE\bfseries Qi Ding}\\[0.16em]
  Ph.D. Candidate in Management Science and Engineering\\
  Institute of Energy, Environment, and Economy, Tsinghua University\\[0.2em]
  \href{mailto:dingq23@mails.tsinghua.edu.cn}{dingq23@mails.tsinghua.edu.cn}
  \enspace$\cdot$\enspace
  \href{https://dingqi01.github.io/}{dingqi01.github.io}
  \enspace$\cdot$\enspace
  \href{https://github.com/dingqi01}{github.com/dingqi01}
\end{center}

\section{Research Profile}
I study how low-carbon transitions move from end-use demand and technology pathways, through spatial infrastructure and supply chains, to hourly power-system operation and resilience. My work combines integrated assessment, optimization, GIS, and data-driven analysis to connect long-term pathways with operational and infrastructure decisions.

\section{Education}
\entry{Tsinghua University}{Ph.D. in Management Science and Engineering, Institute of Energy, Environment, and Economy}{2023--present}
\entry{Tongji University}{B.Eng. in Vehicle Engineering, School of Automotive Studies}{2018--2023}

\section{Peer-Reviewed Publications}
\begin{itemize}
  \item \textbf{Ding, Q.}, Ren, J., Zhang, S., \& Chen, W. (2026). The role of shared autonomous electric vehicles in decarbonizing China's passenger transport sector. \textit{Applied Energy}, 422, 128259. \href{https://doi.org/10.1016/j.apenergy.2026.128259}{[DOI]}
  \item \textbf{Ding, Q.}, \& Chen, W. (2026). Retired Battery Supply and Energy Storage Demand in China: A Provincial-Scale Assessment. \textit{Energy Proceedings}, 65. \href{https://doi.org/10.46855/energy-proceedings-12238}{[DOI]}
  \item Chen, W., Zhang, S., Zhang, Q., Ren, J., \& \textbf{Ding, Q.} (2025). Assessing China's Province-Coordinated Power System Carbon-Neutral Transition Pathway. \textit{Journal of Energy and Climate Change}, 1(1), 1--15. \href{https://doi.org/10.3724/j.issn.2097-4981.JECC-2024-0022}{[DOI]}
\end{itemize}

\section{Preprint and Manuscript}
\begin{itemize}
  \item Pathak, A., Vishwanathan, S. S., Fujimori, S., et al. (including \textbf{Qi Ding}). Commonalities and differences in national pathways toward net-zero emissions. \textit{Research Square} preprint; revision in progress. \href{https://doi.org/10.21203/rs.3.rs-8802998/v1}{[DOI]}
  \item \textbf{Ding, Q.}, \& Chen, W. Aligning Retired EV-Battery Supply with Grid-Storage Demand through Interprovincial Coordination in China. Manuscript under review.
\end{itemize}

\section{Selected Research Projects}
\begin{itemize}
  \project{National Key R\&D Program of China}{CO$_2$-EOR and Long-term Sequestration}{Core member. Techno-economic assessment of CCUS pathways, sequestration benefits, and cost curves.}{Nov 2024--present}
  \project{European Commission, DG CLIMA}{COMMITTED Project}{Research contributor. Multi-model comparison of climate-policy portfolios and transition pathways.}{Sep 2024--present}
  \project{Energy Foundation China}{China Energy System Model Comparison}{Research member. Comparison of model assumptions and decarbonization pathways across energy-system models.}{Sep 2023--present}
  \project{The World Bank Group}{Green China}{Researcher. China TIMES 2.0 analysis of low-carbon power and industry transition pathways.}{Sep 2023--Jun 2024}
  \project{CNPC Economics \& Technology Research Institute}{Transportation-Energy Infrastructure Upgrade}{Research member. Assessment of multi-energy hubs spanning charging, battery swapping, and hydrogen.}{Jul--Dec 2023}
  \project{CATARC}{Commercial-Vehicle Carbon-Neutrality Roadmap 1.0}{Core member at the China Automotive Technology and Research Center. Assessed electrification, hydrogen, and biofuels for the clean-energy-supply workstream.}{Oct 2022--Jul 2023}
\end{itemize}

\newpage

\section{Current Research}
\begin{itemize}
  \currentitem{EV flexibility and power-system resilience under climate stress}{Working paper}{Historically grounded climate-stress events, fixed-capacity power-system replay, smart charging, and vehicle-to-grid operation.}
  \currentitem{Spatial supply chains for hydrogen-derived fuels}{Research in progress}{GIS-informed multi-period optimization of H$_2$, NH$_3$, methanol, and sustainable-aviation-fuel-equivalent production and transport.}
  \currentitem{Charging-infrastructure retrofit planning}{Model development}{Station-adaptive photovoltaic and battery sizing under economic and lifecycle-emissions objectives.}
  \currentitem{Coupling long-term pathways with hourly power-system analysis}{Method development}{Traceable scenario interfaces between China TIMES outputs and spatially and temporally explicit PyPSA-China inputs.}
\end{itemize}

\section{Selected Professional Experience}
\experience{Mianyang Municipal SASAC}{Summer Research Intern}{Contributed to preliminary research for the municipal state-owned-enterprise development plan and to analysis of industrial investment and coordination mechanisms.}{Jul--Aug 2025}
\experience{IM Motors}{Innovation Strategy Intern}{Researched electric-vehicle value chains, technology pathways, and company strategies; synthesized evidence for internal discussions on supply chains and technology investment.}{Sep 2022--Apr 2023}
\experience{NIO}{Project Management Intern, Software Delivery}{Supported production-vehicle software releases by tracking issues, coordinating development and testing interfaces, and contributing to version-review workflows.}{Sep--Dec 2022}

\section{Selected Presentations}
\begin{itemize}
  \item \textbf{Ding, Q.}, et al. Integrating End-Use Flexibility and Climate-Stress Analysis into Sector-Coupled Energy Transition Pathways. Oral presentation, 19th Annual Meeting of the Integrated Assessment Modeling Consortium, 2026.
  \item \textbf{Ding, Q.}, et al. Operational Resilience Value of EV Flexibility under Historically Grounded Climate Stress. Oral presentation, Tsinghua Heli Anbang Joint Doctoral Academic Forum, 2025.
  \item \textbf{Ding, Q.}, et al. Economic-Carbon Misalignment in Photovoltaic-Storage Retrofits for 1,000 Public Charging Stations in Beijing. Oral presentation, 20th National Environmental Doctoral Academic Conference, 2025.
\end{itemize}

\section{Report Contribution}
\begin{itemize}
  \item \textbf{Ding, Q.}, et al. China's Fifth National Climate Change Assessment Report: GHG Mitigation Policy Actions and Outcomes. Supporting analysis for the chapter on mitigation-policy actions, sectoral progress, and carbon-neutral transition pathways, 2025.
\end{itemize}

\section{Models, Methods, and Skills}
\begin{itemize}
  \item \textbf{Transition and power-system models:} China TIMES 2.0, China TIMES-30PE, and PyPSA-China.
  \item \textbf{Decision analysis:} linear and mixed-integer optimization, scenario analysis, and techno-economic assessment.
  \item \textbf{Spatial and data methods:} GIS, network analysis, lifecycle assessment, and data-driven modeling.
  \item \textbf{Research computing:} Python, GAMS, MATLAB, Git, and reproducible model and data workflows.
  \item \textbf{Languages:} English (CET-6: 641) and German (Band 6: Excellent).
\end{itemize}

\section{Selected Honors}
National Scholarship (2021, 2022); Shanghai Outstanding Graduate (2023).

\end{document}
